\subsection{Choosing the right collection (E3)}
\label{sec:routing-audit}
\label{sec:exp-e3}
The native comparison in Section~\ref{sec:e10} did not exercise collection
selection: the agent made no catalogue calls, and at least one search in
nine questions lacked an effective five-collection restriction. Figure~\ref{fig:case-native}a shows the
agent going straight to search. To test that step, we compared two
catalogue-card variants with a scripted selector. A0 supplied the name and a
description derived from the collection's contents; A3 was intended to
supply the name alone. With
descriptions, the selector chose the labelled collection on all 72
questions, against 37 under A3 in the frozen run (Figure~\ref{fig:routing}).
The direction persisted on the two reruns. A3 could fall back to
descriptions or search outside the evaluated collections; those paths were
not logged, so the contrast does not isolate the effect of descriptions.

Reaching the right collection is only a first step. Even when A0 identified
it correctly, A0 returned the target string on only 37 questions. Searching
in the right place does not ensure that the right passage reaches the reader.

\begin{figure}[!htb]
\centering
\begin{minipage}{\columnwidth}
\centering
\includegraphics[width=0.96\linewidth]{fig_opt_route}
\caption{Correct collection selections for the two route-card variants.
All three runs use the same 72 questions. A3 is nominally name-only.}
\label{fig:routing}
\end{minipage}
\end{figure}

\subsection{Getting the target text into the evidence pack (E4)}
\label{sec:ladder-opt}
\label{sec:exp-e4}
We next varied how evidence is assembled once a collection is chosen.
Table~\ref{tab:ablation-configs} defines the six configurations. They differ
in search scope (one collection selected by name and description, or all
ten), in whether documents are read and gated, and in how the pack is
ordered. Routed vector search requests up to 30 candidates at a cosine
threshold of 0.35 and retries at 0.30 only if it returns nothing; the
unrouted arms search each collection at 0.35. Some reference lookups could
use the gold domain, but it did not determine the route-card choice
(Section~\ref{sec:threats}).

\begin{table}[!ht]
\centering
\small
\setlength{\tabcolsep}{3pt}
\renewcommand{\arraystretch}{1.16}
\caption{Search scope and evidence construction for the six E4 configurations.
Limits are in characters, not model tokens.}
\label{tab:ablation-configs}
\begin{tabular}{@{}l>{\raggedright\arraybackslash}p{0.23\columnwidth}>{\raggedright\arraybackslash}p{\dimexpr0.77\columnwidth-31pt\relax}@{}}
\toprule
Arm & Search scope & Candidate selection and evidence\\
\midrule
A0 & One collection, selected by name and description & Retrieve vector hits. The reference-mode gate reads up to six distinct documents, with a nominal 40{,}000-character window per document.\\
A0f & Same route as A0 & Use the same candidates with an inline gate. Allow up to 250{,}000 characters per candidate and in the gate pack.\\
A0h & Same route as A0 & Keep the A0f gate and add a second vector search. Place up to eight chunk windows first, then add gate-retained text up to a final limit of 60{,}000 characters.\\
PD & All ten collections & Keep the six highest-scoring distinct collection/path pairs and use the A0f inline gate.\\
A1b & All ten collections & Merge vector hits and use the A0 reference-mode gate, including document reads.\\
A2 & All ten collections & Merge the same hits as A1b and concatenate their snippets. Do not read documents or run a gate.\\
\bottomrule
\end{tabular}
\end{table}

In the frozen run, A0h retrieved the annotated string on 52 of 72 questions,
up from 37 with A0. The gain came from questions that A0 had missed: A0h
recovered 15 of them without losing any A0 hit (exact McNemar
$p=6.1\times10^{-5}$). Both same-question reruns preserved the gain. A0h
also led the unrouted vector union A1b, but its difference from the unrouted
per-document configuration PD was not significant
(Figure~\ref{fig:content}).

These contrasts compare configurations rather than single components. A0h
changes both what the system retrieves and what it puts first in the pack, so
its comparison with A0f is not a test of window size alone. A2 removes
document reads as well as the gate, so it does not isolate the gate's
effect. The configurations also differ in retrieval-text volume. The trace
counted 309~kB of text per question for A0h, against 31.4~kB for A0 and
24.3~MB for the mechanical read-all scan, which found every target. For A0h,
the count sums candidate text fetched for the gate, the gate's output and
the selected chunk windows, so text can be counted more than once. The final
A0h pack has a separate 60{,}000-character cap; these byte totals are neither
final pack sizes nor generator input sizes.

\begin{figure*}[!tb]
\centering
\includegraphics[width=\textwidth]{fig_opt_content}
\caption{Target-string containment on 72 in-domain questions. (a) Circles
and Wilson 95\% intervals show the frozen run; open triangles show the
same-question follow-up. The dashed line marks frozen A0. The columns give
hit counts out of 72; a dash indicates that the read-all follow-up was
incomplete. (b) For each frozen-run contrast, the symbols count questions
on which only the first or only the second configuration hit. McNemar
$p$ values are exact, two-sided and unadjusted.}
\label{fig:content}
\end{figure*}

\subsection{Storing and finding operational lessons (E5)}
\label{sec:e3}
\label{sec:exp-e5}
HMRS stores operational lessons as searchable documents and orders them by
credibility tier. Only parts of this workflow could be verified. A smoke test
produced candidate drafts but found no global hits. A separate case report
placed a stored lesson first in a search; no answer trace was retained. The
approval check found an empty draft pool, and the heuristic dry run did not
meet its quality criteria. These checks illustrate parts of the workflow;
they do not show reliable recovery or automatic approval of lessons in
deployment.
